
%\documentclass[journal]{IEEEtran}
\documentclass[journal]{igtesymp}
% *** GRAPHICS RELATED PACKAGES ***
%
\ifCLASSINFOpdf
  % \usepackage[pdftex]{graphicx}
  % declare the path(s) where your graphic files are
  % \graphicspath{{../pdf/}{../jpeg/}}
  % and their extensions so you won't have to specify these with
  % every instance of \includegraphics
  % \DeclareGraphicsExtensions{.pdf,.jpeg,.png}
\else
  % or other class option (dvipsone, dvipdf, if not using dvips). graphicx
  % will default to the driver specified in the system graphics.cfg if no
  % driver is specified.
  % \usepackage[dvips]{graphicx}
  % declare the path(s) where your graphic files are
  % \graphicspath{{../eps/}}
  % and their extensions so you won't have to specify these with
  % every instance of \includegraphics
  % \DeclareGraphicsExtensions{.eps}
\fi
% graphicx was written by David Carlisle and Sebastian Rahtz. It is
% required if you want graphics, photos, etc. graphicx.sty is already
% installed on most LaTeX systems. The latest version and documentation can
% be obtained at: 
% http://www.ctan.org/tex-archive/macros/latex/required/graphics/
% Another good source of documentation is "Using Imported Graphics in
% LaTeX2e" by Keith Reckdahl which can be found as epslatex.ps or
% epslatex.pdf at: http://www.ctan.org/tex-archive/info/
%
% latex, and pdflatex in dvi mode, support graphics in encapsulated
% postscript (.eps) format. pdflatex in pdf mode supports graphics
% in .pdf, .jpeg, .png and .mps (metapost) formats. Users should ensure
% that all non-photo figures use a vector format (.eps, .pdf, .mps) and
% not a bitmapped formats (.jpeg, .png). IEEE frowns on bitmapped formats
% which can result in "jaggedy"/blurry rendering of lines and letters as
% well as large increases in file sizes.
%
% You can find documentation about the pdfTeX application at:
% http://www.tug.org/applications/pdftex


% *** SUBFIGURE PACKAGES ***
\usepackage[tight,footnotesize]{subfigure}
% subfigure.sty was written by Steven Douglas Cochran. This package makes it
% easy to put subfigures in your figures. e.g., "Figure 1a and 1b". For IEEE
% work, it is a good idea to load it with the tight package option to reduce
% the amount of white space around the subfigures. subfigure.sty is already
% installed on most LaTeX systems. The latest version and documentation can
% be obtained at:
% http://www.ctan.org/tex-archive/obsolete/macros/latex/contrib/subfigure/
% subfigure.sty has been superceeded by subfig.sty.



\usepackage{graphicx}
%\usepackage{german}
%\usepackage{umlaute}


% *** Do not adjust lengths that control margins, column widths, etc. ***
% *** Do not use packages that alter fonts (such as pslatex).         ***
% There should be no need to do such things with IEEEtran.cls V1.6 and later.
% (Unless specifically asked to do so by the journal or conference you plan
% to submit to, of course. )



% correct bad hyphenation here
\hyphenation{op-tical net-works semi-conduc-tor}


\begin{document}
%
% paper title
% can use linebreaks \\ within to get better formatting as desired
\title{Volume Integral Method for Cauer Ladder Network Synthesis: \\
       Double-Double Verification on a Conducting Sphere}
%
\author{K.~Sugahara\IEEEauthorrefmark{1},  H.~Nagamine\IEEEauthorrefmark{2},
and Y.~Hane\IEEEauthorrefmark{3}
 \\ \vspace{0.3cm} \normalsize{
\IEEEauthorblockA{\IEEEauthorrefmark{1}Faculty of Science and Engineering, Kindai University, 3-4-1 Kowakae, Higashi-Osaka, Osaka 577-8502, Japan\\
 \IEEEauthorrefmark{2}Faculty of Engineering, Gifu University, 1-1 Yanagido, Gifu 501-1193, Japan\\
 \IEEEauthorrefmark{3}Faculty of Science and Engineering, Toyo University, Kawagoe, Saitama, Japan\\
 E-mail: ksugahar@kindai.ac.jp}}%
}


% note the % following the last \IEEEmembership and also \thanks - 
% these prevent an unwanted space from occurring between the last author name
% and the end of the author line. i.e., if you had this:
% 
% \author{....lastname \thanks{...} \thanks{...} }
%                     ^------------^------------^----Do not want these spaces!
%
% a space would be appended to the last name and could cause every name on that
% line to be shifted left slightly. This is one of those "LaTeX things". For
% instance, "\textbf{A} \textbf{B}" will typeset as "A B" not "AB". To get
% "AB" then you have to do: "\textbf{A}\textbf{B}"
% \thanks is no different in this regard, so shield the last } of each \thanks
% that ends a line with a % and do not let a space in before the next \thanks.
% Spaces after \IEEEmembership other than the last one are OK (and needed) as
% you are supposed to have spaces between the names. For what it is worth,
% this is a minor point as most people would not even notice if the said evil
% space somehow managed to creep in.




% use for special paper notices
%\IEEEspecialpapernotice{(Invited Paper)}




\IEEEaftertitletext{\vspace{-1cm}\noindent\begin{abstract}
The Cauer Ladder Network (CLN) representation of an eddy-current device is essential for circuit-level loss models in power electronics and electric machines.
The current de-facto extraction is FEM Kameari accumulation, but the underlying Krylov-subspace iteration loses M-orthogonality at higher Cauer stages, and finite-precision arithmetic ultimately limits the reliable number of rungs that can be extracted.
We propose a Volume Integral Method (VIM) extraction in which the conducting body alone is discretised, the diffusion eigenproblem $K v_n = \lambda_n M v_n$ is solved once, and the Kameari-convention Cauer-I rungs are obtained from the Foster amplitudes by Hankel-Pad\'e continued-fraction inversion of the moments $\alpha_n = (-1)^n \sum_k g_k^2 \tau_k^n$.
The formulation is verified on a copper sphere under uniform $B_z$, where Stoll's Bessel-series analytical Cauer ladder serves as ground truth.
A pure-Python/CuPy {\em double-double} ($\sim 32$ digit) implementation of the VIM-CLN pipeline is then introduced, in which the singular axisymmetric Green's-function kernel is evaluated via {\tt mpmath} elliptic integrals and the generalized eigenproblem is solved via Cholesky-based {\tt mpmath} arithmetic.
Combined with mpmath verified-interval Hankel-Pad\'e continued-fraction inversion, the DD pipeline reaches the machine DD precision floor $K_{\text{lo}}/K_{\text{hi}}\!=\!1.1\!\times\!10^{-16}$ and reproduces Stoll's leading rungs $R_0, L_1$ to 0.1\,\% / 1.8\,\% at the demonstrated $540$-cell scale, with verified-interval relative width below $10^{-30}$ across all extracted stages.
The same DD pipeline applies to hex / axisymmetric VIM uniformly, providing the first GPU-accessible verified-arithmetic CLN extractor.
\end{abstract}
\noindent\begin{keywords}
Cauer Ladder Network, Eddy Current, Foster Spectrum, Volume Integral Method.
\end{keywords}\vspace{\baselineskip}}


\maketitle
\thispagestyle{empty}\pagestyle{empty}





% For peer review papers, you can put extra information on the cover
% page as needed:
% \ifCLASSOPTIONpeerreview
% \begin{center} \bfseries EDICS Category: 3-BBND \end{center}
% \fi
%
% For peerreview papers, this IEEEtran command inserts a page break and
% creates the second title. It will be ignored for other modes.
\IEEEpeerreviewmaketitle

\section{Introduction}
Eddy-current loss models are essential for circuit-level design of high-frequency transformers, SiC/GaN power-electronic modules, and high-speed electric machines.
The Cauer Ladder Network (CLN) representation reduces the eddy-current admittance to a finite ladder of resistances $R_n$ and inductances $L_n$, suitable for direct embedding in circuit simulators \cite{Sugahara_TEAM28, Hiruma_2020}.
The current de-facto extractor is the FEM Kameari accumulation iteration; however, this scheme suffers from numerical loss of M-orthogonality in the underlying Krylov subspace, and Hankel-Pad\'e continued-fraction inversion conditioning amplifies floating-point input noise by $\sim 300\times$ per stage so that FP64 arithmetic is restricted to roughly four to five reliable Cauer rungs \cite{Nagamine_2026}.

We propose a {\it Volume Integral Method} (VIM) approach to CLN synthesis.
The conducting body is discretised with HDiv divergence-free volume basis functions for solid hex elements or piecewise-constant $J_\phi$ axisymmetric cells; the diffusion eigenproblem
$K v_n = \lambda_n M v_n$
is solved once, and the Kameari-convention Cauer-I rungs are obtained from the Foster amplitudes by Hankel-Pad\'e continued-fraction inversion.
Because eigenvectors are generated all at once rather than iteratively, the M-orthogonality loss issue is avoided in principle.
We then introduce a {\em double-double} ($\sim 32$ digit) GPU pipeline of the same VIM-CLN extractor and show that the FP64 precision wall can be pushed back by several Cauer stages on a benchmark conducting sphere, where Stoll's analytical Bessel-series solution serves as ground truth \cite{Stoll1974}.

\section{VIM-based CLN Formulation}

For the conductor occupying $\Omega_c$ with conductivity $\sigma$, we expand the volume current density in basis $\phi_i$ adapted to the geometry: HDiv divergence-free hex basis for solid conductors, or piecewise-constant $J_\phi$ on a meridian $(r,z)$ mesh for axisymmetric ones.
The energy matrices are
\begin{equation}
K_{ij} = \frac{\mu_0}{4\pi}\!\!\int_{\Omega_c}\!\int_{\Omega_c} \frac{\phi_i(\mathbf{r}) \cdot \phi_j(\mathbf{r}')}{|\mathbf{r}-\mathbf{r}'|}\,d\mathbf{r}\,d\mathbf{r}', \label{eq:K}
\end{equation}
\begin{equation}
M_{ij} = \int_{\Omega_c} \phi_i \cdot \phi_j \,d\mathbf{r}. \label{eq:M}
\end{equation}
The 3D $1/r$ kernel singularity in (\ref{eq:K}) is integrated with Spherical Duffy quadrature; in the axisymmetric case, the singular $\phi$-integrated Green's function reduces to complete elliptic integrals $K(m), E(m)$ that are evaluated to arbitrary precision via {\tt mpmath}.
The generalized eigenproblem
$K v_n = \lambda_n M v_n$
gives Foster time constants $\tau_n = \sigma \lambda_n$ and modes $v_n$ orthonormal in $M$.
The applied field $\mathbf{A}_{\text{ext}}$ projects onto each mode through $g_n = v_n^\top b$ with $b_i = \int_{\Omega_c} \phi_i \cdot \mathbf{A}_{\text{ext}}\,d\mathbf{r}$.

We adopt the {\it Kameari $\chi$-Foster convention}, in which the Cauer-I ladder is extracted from the Hankel-Pad\'e moments
\begin{equation}
\alpha_n = (-1)^n \sum_k g_k^2 \, \tau_k^n. \label{eq:alpha}
\end{equation}
Note that the literature also reports a {\it Hiruma Krylov-Pad\'e impedance convention} \cite{Hiruma_2020} which uses the moments $\sum_k g_k^2 \tau_k^{n+1}$; the two conventions yield systematically different Cauer ladders ($\sim 5$--$12\,\%$ for the leading rung).
The Kameari convention is chosen here because it coincides with the $\chi$-Foster susceptibility extracted from the Stoll closed-form solution.
Cauer-I rungs $R_{2k}, L_{2k+1}$ are obtained by continued-fraction inversion of $\{\alpha_n\}$.
The continued-fraction inversion conditioning is quantified by {\em mpmath verified-interval arithmetic} \cite{Nagamine_2026}: rungs whose relative interval width exceeds 1\,\% are flagged as numerically unreliable. With FP64 input moments this threshold is crossed at stage 4--5 across all tested geometries.

\section{Sphere Verification}

We test the VIM-CLN extraction on a copper sphere ($a=10$\,mm, $\sigma=5.8\times10^7$\,S/m) under a uniform $B_z=1$\,T applied field, for which Stoll \cite{Stoll1974} gives the analytical $\chi$-Foster representation
$\chi(s) = -1 + \sum_n \frac{6/(n^2\pi^2)}{1 + s\,\mu_0\sigma a^2/(n^2\pi^2)}$.
Hankel-Pad\'e inversion of (\ref{eq:alpha}) at 240-digit Mathematica arithmetic produces 40 reliable Cauer-I rungs as the analytical reference.

The VIM with HDiv order 3 ($n_{\text{dofs}}=81$) and Spherical Duffy quadrature in FP64 reproduces this analytical Cauer ladder to machine precision through stage 4, confirming both the formulation and the convention identification (Table \ref{tab:sphere}, columns labelled ``VIM FP64'').

\begin{table}[h]
\caption{Sphere Cauer-I rung $\tau_{\text{pair}}=L/R$ ($\mu$s): VIM FP64 vs.\ DD pipeline (axisym 270 / 540 cells) vs.\ Stoll analytical (Kameari convention). DD 540 cells extends the reliable Cauer ladder by one stage over the 270-cell baseline (DD 270 stage 5 is already inflated; DD 540 keeps stage 5 finite and shifts the sign-flip breakdown to stage 6).}\label{tab:sphere}
\begin{center}\footnotesize
\begin{tabular}{|c|r|r|r|r|}
\hline
$k$ & Stoll [$\mu$s] & VIM FP64 [$\mu$s] & DD 270 [$\mu$s] & DD 540 [$\mu$s] \\\hline
0 & 694.142 & 694.142 & 720.36 & 706.75 \\\hline
1 & 154.604 & 154.605 & 179.79 & 166.80 \\\hline
2 & 64.075 & 64.075 & 91.59 & 77.79 \\\hline
3 & 34.474 & 34.475 & 62.81 & 49.34 \\\hline
4 & 21.400 & 21.402 & 49.38 & 36.56 \\\hline
5 & --- & --- & 164.49$^\dagger$ & 28.93 \\\hline
\end{tabular}
\end{center}
{\footnotesize $^\dagger$ Already inflated at the 270-cell mesh due to Foster-spectrum truncation; the DD 540 cells run keeps stage 5 within an order of magnitude of analytical and pushes the sign-flip breakdown to stage 6.}
\end{table}

\section{Double-Double GPU Pipeline}

To push the Hankel-Pad\'e input-precision wall back, we re-implemented the entire VIM-CLN pipeline in {\em double-double} (DD, $\sim 32$ decimal digit) arithmetic in pure Python/CuPy: DD primitives (Bailey-Hida-Li), mpmath Newton-refined Gauss-Legendre nodes at 40-digit, DD HDiv basis evaluation for hex elements, mpmath complete elliptic-integral evaluation of the axisymmetric Green's function, Cholesky-based generalised eigh in {\tt mpmath} arithmetic, and verified-interval Hankel-Pad\'e inversion.
The DD precision floor $K_{\text{lo}}/K_{\text{hi}}=1.1\!\times\!10^{-16}$ is reached, confirming that all $\sim 14$ trailing decimal digits below FP64 are correctly captured.
For the sphere at $N_\rho\!\times\!N_\theta\!=\!20\!\times\!27\!=\!540$ axisymmetric cells the DD pipeline matches Stoll's leading $R_0$ to $0.1\,\%$ and the leading inductance $L_1$ to $\sim\!1.8\,\%$ (compared to $0.2\,\%/4.0\,\%$ at the 270-cell baseline -- doubling the cells roughly halves the gap), while the verified-interval relative width remains below $10^{-30}$ across all extracted stages -- so any gap to the analytical result is attributable to Foster-spectrum discretisation, not to floating-point conditioning.
The same DD pipeline operates uniformly on hex (3D solid) and axisymmetric VIM discretisations; the GPU CuPy port of the hex K-assembly delivers a $4300\times$ speed-up over the CPU C++ implementation, enabling DD HDiv-basis hex production runs at order $p=4$ ($n_{\text{dofs}}=176$) in seconds, while the axisymmetric K-assembly is multiprocessed across CPU cores.

\section{Conclusion}

A Volume Integral Method approach to Cauer Ladder Network synthesis was proposed and verified against Stoll's analytical Bessel solution for a conducting sphere with machine precision through five Cauer rungs in FP64.
The Kameari $\chi$-Foster convention was identified as the natural choice for VIM extraction, distinct from the Hiruma Krylov-Pad\'e impedance convention reported in earlier literature.
A double-double precision pipeline was then implemented end-to-end in pure Python/CuPy, providing the first GPU-accessible verified-arithmetic CLN extractor and a foundation on which the precision-wall analysis of Nagamine's verified-interval Hankel-Pad\'e arithmetic can be applied directly to electromagnetic eddy-current problems.

\begin{thebibliography}{1}
\bibitem{Sugahara_TEAM28} K.~Sugahara, ``Eddy current Cauer ladder network extraction for TEAM problem 28,'' Compumag 2017.
\bibitem{Hiruma_2020} S.~Hiruma and H.~Igarashi, ``Three-term recurrence relation for Cauer ladder synthesis of eddy-current problems,'' IEEE Trans.\ Magn., vol.\ 56, no.\ 3, 2020.
\bibitem{Stoll1974} R.~L.~Stoll, {\it The Analysis of Eddy Currents}. Oxford: Clarendon, 1974.
\bibitem{Nagamine_2026} M.~Nagamine et al., ``Verified numerical computation of the Cauer network representation of a square prism conductor,'' submitted, 2026.
\bibitem{radia} K.~Sugahara, ``Radia: Electromagnetic Simulation Framework for Magnetic Levitation Systems,'' GitHub repository, https://github.com/ksugahar/Radia, accessed May~2026.
\end{thebibliography}

% that's all folks
\end{document}
